\section{Contextual and Specialized Setups}

\subsection{Spring trades: stretched impulse mean reversion}

The ``spring'' setup describes a market that moves too far too fast in one direction. The key
idea is not simply overbought / oversold language. The key idea is that a one-directional
impulse through a thin book can stretch the local market so much that even small opposing
passive interest becomes meaningful.

The clean version is: one strong impulse, little pullback, obvious stretching, then visible
resistance appears near the extreme. The trader does not need enormous opposing size because
the stretch itself supplies the vulnerability. In the transcript example, a small cluster of
passive orders becomes enough because the market is already ``wound up.''

\paragraph{Algorithmic translation.}
Measure one-way displacement over short windows, lack of intermediate retracements, local book
emptiness, and first opposing liquidity near the extreme. This setup is really a context-
dependent bounce: weak passive resistance can become strong \emph{because} the path to it was so
imbalanced.

\subsection{Plank or limit-state trades}

The plank lesson is explicit about danger. When a market reaches its daily limit state, price
behavior becomes non-linear. Expansion timing is uncertain. When the limit expands, price often
shoots in the direction of the plank first, but the environment is described as ``inadequate''
or structurally abnormal. That is a warning, not an invitation.

This setup belongs in the catalog because it matters operationally, but it should be treated as
a specialist regime. A future research engine should classify limit-state behavior as a regime
change with separate rules rather than as a standard continuation or fade setup.

\paragraph{Algorithmic translation.}
Detect proximity to exchange-imposed daily limit boundaries, time spent pinned, repeated queue
stacking near the boundary, and first post-expansion impulse. Confidence should remain low until
the platform has actual live data on enough such episodes.

\subsection{News reaction setups}

The news lesson gives one of the most useful research principles in the entire knowledge base:
the trader does not try to interpret whether the news is ``good'' or ``bad'' in an abstract
sense. The trader watches how the market reacts. News is treated as a warning that a strong move
may appear. The actual trade still comes from order flow.

This has several implications:

\begin{itemize}[leftmargin=1.4em]
  \item News is a \term{context activator}, not a complete trading signal.
  \item Reaction quality matters more than headline semantics.
  \item Different instruments can express the same news through different leader channels.
  \item The right question is not ``what should happen?'' but ``what are participants doing now?''
\end{itemize}

\paragraph{Algorithmic translation.}
Use scheduled-news calendars as binary context flags, then evaluate post-event volatility burst,
depth thinning, trade intensity, and the first accepted direction. Event studies around fixed
news windows are more realistic than trying to infer narrative sentiment from the source
material alone.

\subsection{Guide-instrument, spot-futures, and leader-follower logic}

One transcript makes a strong claim that several related instruments can trade as one logical
instrument. The clearest example is spot, tomorrow settlement, and futures in currency markets.
If a meaningful iceberg or density appears in the leader, the followers may stall until it is
resolved. This is a concrete form of leader-follower logic rather than a vague correlation
claim.

For research purposes, this means cross-instrument event linking is essential. A density event
in one venue or settlement type can explain the delayed behavior of a related contract. This is
especially relevant for the Russian market, where related contracts and settlement variants can
interact tightly.

\paragraph{Algorithmic translation.}
Record synchronized books and trades for related instruments. Detect whether an event in the
leader precedes stalling, breakout, or bounce in the follower. Measure lag, consistency, and
whether the relationship is stable only at certain times of day.

\subsection{PoC, cluster, and level-based extensions}

The Professional Scalping book adds a broader family of setups around historical volume
clusters and Point of Control. The key conceptual phrase is excellent: PoC can be treated as
``densities from the past.'' This generalizes the DOM idea. Some liquidity is visible \emph{now};
other relevant interest is inferred from where the market previously did the most business.

That gives a hierarchy of levels:

\begin{itemize}[leftmargin=1.4em]
  \item obvious chart levels,
  \item current visible DOM density,
  \item historical cluster or PoC concentration,
  \item smart-money style zones inferred from large-participant behavior.
\end{itemize}

The practical lesson is not that every PoC should be traded. The practical lesson is that
current order-book events gain meaning when they happen at a historically important price.

\subsection{Smart Money framing}

The book's Smart Money framing is conceptually useful but operationally fuzzy. It argues that
large institutions and major prop participants move markets and that the trader's job is to read
their intentions. This overlaps with iceberg logic, robot logic, and leader-follower logic. In
practice, Smart Money should be interpreted as a \term{meta-description} of several smaller
observable mechanisms rather than as one direct trigger.

\subsection{Spread capture and volatility harvesting}

Even though the dedicated local transcript file is empty, the broader materials imply a class of
trades focused less on major directional bets and more on short-lived microstructure dislocation:
wide spread, local inefficiency, temporary imbalance, and rapid normalization. This is likely
what later automation could exploit best in liquid markets, especially in crypto.

However, this remains weakly specified in the current repository. It should therefore be treated
as a future research branch rather than a fully formalized setup today.
